\chapter{Supplements, labels, and interactions}
\section{Learning objectives}
Evaluate the quality of a supplement claim; read a product label; identify pharmacokinetic and pharmacodynamic interactions; and communicate a safe monitoring and deprescribing plan.

\section{Reading a label}
Check serving size, chemical form, elemental mineral amount, percent Daily Value, and whether a serving contains multiple products you already take. ``Natural'' does not mean risk-free. Gummies may resemble candy, powders may be concentrated, and proprietary blends can hide individual amounts. Choose products with credible independent quality testing when available; testing does not prove that a product is necessary or effective.

\section{Common interaction patterns}
Interactions may be pharmacokinetic, changing absorption, metabolism, transport, or excretion, or pharmacodynamic, producing additive effects such as bleeding, sedation, hypertension, hypoglycemia, or arrhythmia. Timing separation can help with some binding interactions but does not remove systemic interactions.
Calcium, iron, magnesium, and zinc can reduce absorption of some medicines, including certain antibiotics and thyroid hormone, when taken together. Vitamin K can alter warfarin management. St. John's wort is not a vitamin and can reduce levels of many medicines. High-dose antioxidants may be unhelpful or harmful around some cancer treatments. Biotin can interfere with laboratory assays. Always provide clinicians and pharmacists a complete list, including teas, powders, and sports products.

\section{Who should be especially cautious}
Discuss supplementation before use if pregnant or breastfeeding, under 18, preparing for surgery, living with kidney or liver disease, having a history of kidney stones or iron overload, or taking anticoagulants, thyroid medicine, seizure medicine, chemotherapy, or medicines for blood pressure and diabetes.

\section{A simple decision test}
Ask: What problem am I trying to solve? Is there evidence that this nutrient helps that problem? Can food address it? What is the total dose from every product? What is the plan to stop or reassess? If these questions cannot be answered, pause and seek professional advice.
